\chapter*{Resumo}
\label{ch:resumo}

A crescente geração de \glsauto{fw} e o excedente de \glsauto{cgl} proveniente da produção de biodiesel representam dois fluxos significativos de resíduos orgânicos que requerem soluções de gestão sustentáveis. Este estudo avaliou a viabilidade técnica de um sistema de codigestão anaeróbia de uma única etapa, tratando \glsauto{fw} e \glsauto{cgl} como co-substratos, com o objetivo de verificar a produção de biogás e, simultaneamente, proporcionar uma via de valorização para o \glsauto{cgl}. A investigação foi conduzida num \glsauto{cstr} de \qty{1}{\litre}, operado em condições mesofílicas (\qty{35}{\celsius}), com um \glsauto{hrt} de \num{30} dias. \Glsentrylongauto{fw} foram recolhidos numa cantina universitária, e o \glsauto{cgl} foi obtido a partir de resíduos da produção de biodiesel; lamas de um reator anaeróbio de uma \glsauto{wwtp} municipal serviram como inóculo. O reator foi operado em fases sequenciais: 1–um período inicial de adaptação; 2–codigestão de \glsauto{fw} e \glsauto{cgl}; 3–alimentação exclusiva com \glsauto{cgl}; e 4–um período pós-alimentação. A produção de biogás foi monitorizada continuamente utilizando \glsauto{ampts}, enquanto a estabilidade do processo foi avaliada através de medições de pH, \glsauto{alk} total, \glsauto{vfa} e respectiva \glsauto{vfa_alk}. A caracterização dos substratos incluiu análises de \glsauto{ts}, \glsauto{vs}, massa volúmica, \glsauto{tc} e \glsauto{tn}. Os resultados demonstraram que a Fase 2 (codigestão de \glsauto{fw} e \glsauto{cgl}) originou as taxas de produção de \glsauto{ch4} e os volumes acumulados mais elevados, com picos de fluxo imediatamente após as alimentações, indicando uma rápida conversão do substrato. Em contraste, a Fase 3 (alimentação exclusiva com \glsauto{cgl}) resultou numa produção reduzida de biogás, apesar da atividade metanogénica continuada. A \glsauto{alk} total manteve-se elevada após a suplementação com bicarbonato, mantendo o pH em torno de \num{7} durante a maior parte da experiência; as concentrações de \glsauto{vfa} flutuaram entre \qtyrange{3000}{5000}{\milli\gram\CHCOOH\per\litre}, com \glsauto{vfa_alk} de \numrange{0,3}{1,2}. Estes resultados apoiam a hipótese de que o \glsauto{cgl} pode ser útil para a produção de biogás quando codigerido com \glsauto{fw} ou isoladamente. O estudo demonstra a viabilidade técnica da codigestão anaeróbia como estratégia de valorização do \glsauto{cgl} no âmbito de uma bioeconomia circular, embora seja necessária investigação adicional para otimizar as taxas de carga de \glsauto{cgl} e desenvolver estratégias de alimentação robustas para aplicação industrial.

\vspace{1em}
\noindent\textbf{Palavras-chave:} \palavraschaves.
